\documentclass[aspectratio=169]{beamer}
\usetheme{Madrid}
\usecolortheme{default}
\usepackage[T1]{fontenc}
\usepackage[utf8]{inputenc}
\usepackage[french]{babel}
\usepackage{booktabs}
\usepackage{graphicx}
\usepackage{array}
\usepackage{hyperref}
\usepackage{amsmath}

\title{Simulation probabiliste de politiques de maintenance}
\subtitle{Projet M2D 2025--2026}
\author{Équipe projet}
\date{\today}

\begin{document}

\begin{frame}
  \titlepage
\end{frame}

\begin{frame}{Plan de la présentation}
\tableofcontents
\end{frame}

\section{Problème et objectifs}

\begin{frame}{Problématique et objectif}
\begin{itemize}
  \item On étudie une flotte de systèmes possédant un composant critique qui se dégrade avec l'usage.
  \item L'objectif est de comparer plusieurs politiques de maintenance sur trois axes :
  \begin{itemize}
    \item nombre de pannes,
    \item coût global,
    \item disponibilité.
  \end{itemize}
  \item Le livrable principal est un simulateur stochastique paramétrable.
\end{itemize}

\vspace{0.3cm}
\textbf{Idée centrale :} construire un modèle simple, cohérent et exploitable plutôt qu'un modèle artificiellement complexe.
\end{frame}

\begin{frame}{Questions auxquelles on veut répondre}
\begin{itemize}
  \item Quelle politique limite le mieux le risque de panne ?
  \item Quelle politique minimise le coût total sur l'horizon étudié ?
  \item Comment évolue l'état global de la flotte au cours du temps ?
  \item Dans quelle mesure les conclusions sont-elles robustes à l'aléatoire ?
\end{itemize}

\vspace{0.25cm}
Nous avons donc construit un simulateur probabiliste, puis comparé plusieurs politiques de maintenance et analysé les résultats par Monte Carlo.
\end{frame}

\section{Modélisation}

\begin{frame}{Hypothèses de modélisation}
\begin{itemize}
  \item Tous les systèmes sont supposés identiques.
  \item Chaque système possède un seul composant critique.
  \item Le composant se dégrade progressivement sous l'effet de l'usage et du vieillissement.
  \item Une inspection n'observe pas l'usure exacte, mais un score bruité.
  \item Une réparation est partielle ; un remplacement remet le composant à neuf.
\end{itemize}

\vspace{0.2cm}
Ces hypothèses rendent le modèle lisible, stochastique et cohérent avec le sujet.
\end{frame}

\begin{frame}{Hypothèses et lois probabilistes}
\begin{columns}[T]
\column{0.56\textwidth}
\begin{itemize}
  \item Usage journalier : loi \textbf{log-normale}
  \item Incrément d'usure : loi \textbf{log-normale} par défaut, variante \textbf{Weibull} possible
  \item Inspection : score bruité
  $$score = usure + \varepsilon, \quad \varepsilon \sim \mathcal{N}(0,\sigma^2)$$
  \item Panne : risque journalier croissant avec l'usure et l'âge
  $$p_t = 1 - e^{-\lambda_t}, \quad \lambda_t = \lambda_0 e^{a\,usure + b\,age}$$
\end{itemize}

\column{0.40\textwidth}
\begin{block}{Pourquoi ce choix ?}
\begin{itemize}
  \item éviter un modèle trop déterministe,
  \item garantir des variables positives,
  \item rester lisible et justifiable.
\end{itemize}
\end{block}
\end{columns}
\end{frame}

\begin{frame}{Chaîne logique du simulateur}
\begin{enumerate}
  \item Tirage de l'usage journalier
  \item Mise à jour de l'usure
  \item Inspection éventuelle
  \item Calcul du risque journalier de panne
  \item Tirage de panne ou non
  \item Maintenance éventuelle
  \item Enregistrement des données dans les CSV
\end{enumerate}

\vspace{0.25cm}
On obtient ainsi une simulation \textbf{événementielle simple} mais suffisamment riche pour comparer des politiques.
\end{frame}

\section{Politiques comparées}

\begin{frame}{Politiques comparées}
\centering
\begin{tabular}{>{\raggedright\arraybackslash}p{0.28\textwidth} >{\raggedright\arraybackslash}p{0.58\textwidth}}
\toprule
\textbf{Politique} & \textbf{Principe} \\
\midrule
Corrective pure & Aucune inspection, intervention uniquement après panne \\
Conditionnelle 30 jours & Inspection tous les 30 jours, maintenance sur alerte \\
Systématique 200 jours & Remplacement automatique à 200 jours \\
Mixte 30 jours + 200 jours & Inspection périodique + remplacement systématique \\
\bottomrule
\end{tabular}

\vspace{0.4cm}
Ces quatre politiques couvrent les grandes familles demandées dans le sujet : corrective, conditionnelle, préventive systématique et mixte.
\end{frame}

\begin{frame}{Lecture intuitive des politiques}
\begin{itemize}
  \item \textbf{Corrective pure} : on attend la panne.
  \item \textbf{Conditionnelle 30 jours} : on observe régulièrement et on agit si le détecteur alerte.
  \item \textbf{Systématique 200 jours} : on remplace sans attendre l'apparition d'un défaut détecté.
  \item \textbf{Mixte} : on combine observation régulière et borne d'âge maximale.
\end{itemize}

\vspace{0.2cm}
L'intérêt de la comparaison est de comprendre l'arbitrage entre :
\begin{itemize}
  \item coûts de prévention,
  \item coûts de panne,
  \item qualité globale de la flotte.
\end{itemize}
\end{frame}

\section{Résultats principaux}

\begin{frame}{Résultats principaux -- coût total}
\centering
\includegraphics[width=0.92\textwidth]{presentation_figures/couts_politiques.png}

\vspace{0.2cm}
\small
La politique corrective pure est nettement dominée en coût. Les politiques conditionnelle et systématique offrent ici les meilleurs compromis économiques.
\end{frame}

\begin{frame}{Résultats principaux -- pannes et disponibilité}
\centering
\includegraphics[width=0.92\textwidth]{presentation_figures/pannes_disponibilite.png}

\vspace{0.2cm}
\small
La politique mixte réduit davantage les pannes, mais au prix d'une maintenance plus lourde. La corrective pure garde une bonne disponibilité apparente car elle n'arrête les systèmes qu'après panne.
\end{frame}

\begin{frame}{Résultats principaux -- pannes cumulées au cours du temps}
\centering
\includegraphics[width=0.94\textwidth]{presentation_figures/pannes_cumulees_temps.png}

\vspace{0.15cm}
\small
Cette figure est directement liée au sujet : elle montre comment les pannes s'accumulent au cours du temps selon la politique choisie. Elle est facile à commenter et très parlante en soutenance.
\end{frame}

\begin{frame}{État global de la flotte -- systèmes dégradés}
\centering
\includegraphics[width=0.94\textwidth]{presentation_figures/proportion_systemes_degrades.png}

\vspace{0.15cm}
\small
Cette courbe montre la part de la flotte qui a dépassé le seuil de défaut. Elle décrit la santé globale du parc au fil du temps.
\end{frame}

\begin{frame}{État global de la flotte -- usure moyenne}
\centering
\includegraphics[width=0.94\textwidth]{presentation_figures/usure_moyenne_flotte.png}

\vspace{0.15cm}
\small
La politique préventive agit directement sur la dynamique moyenne d'usure. Cette figure est plus proche du mécanisme physique du modèle que la seule disponibilité.
\end{frame}

\begin{frame}{Lecture croisée des résultats}
\begin{itemize}
  \item La politique \textbf{corrective pure} limite les actions préventives mais paie un prix élevé en pannes.
  \item La politique \textbf{conditionnelle} garde une bonne maîtrise du coût tout en s'appuyant sur le détecteur.
  \item La politique \textbf{systématique} reste compétitive en coût grâce à la réduction du risque de panne.
  \item La politique \textbf{mixte} est plus conservatrice sur l'état de la flotte, mais plus coûteuse.
\end{itemize}
\end{frame}

\section{Robustesse statistique}

\begin{frame}{Robustesse statistique : Monte Carlo}
\centering
\includegraphics[width=0.92\textwidth]{presentation_figures/monte_carlo_couts.png}

\vspace{0.2cm}
\small
Plusieurs réplications sont nécessaires pour quantifier la variabilité. Les intervalles de confiance à 95\% confirment la hiérarchie générale entre politiques.
\end{frame}

\begin{frame}{Pourquoi l'analyse Monte Carlo est utile}
\begin{itemize}
  \item Une seule simulation peut être influencée par un tirage favorable ou défavorable.
  \item Les réplications permettent de calculer une moyenne et un intervalle de confiance.
  \item On sépare ainsi ce qui relève de la tendance générale et ce qui relève du hasard.
\end{itemize}

\vspace{0.2cm}
\textbf{Message à retenir :} les conclusions présentées ne reposent pas sur un unique scénario simulé.
\end{frame}

\section{Analyse du détecteur}

\begin{frame}{Qualité du détecteur : métriques type ML}
\begin{columns}[T]
\column{0.52\textwidth}
\centering
\includegraphics[width=\textwidth]{ml_metrics/conditionnelle_30j/roc_curve.png}

\column{0.44\textwidth}
\small
\begin{tabular}{lr}
\toprule
\textbf{Métrique} & \textbf{Valeur} \\
\midrule
Précision & 0.711 \\
Rappel & 0.777 \\
F1-score & 0.742 \\
ROC AUC & 0.919 \\
PR AUC & 0.827 \\
\bottomrule
\end{tabular}

\vspace{0.35cm}
Le détecteur simulé peut être interprété comme un classifieur binaire sur les inspections.
\end{columns}
\end{frame}

\begin{frame}{Comparaison synthétique des scores ML}
\centering
\includegraphics[width=0.86\textwidth]{presentation_figures/roc_auc_policies.png}

\vspace{0.2cm}
\small
Les politiques inspectées disposent d'un détecteur globalement performant. Les métriques ML complètent l'analyse économique en évaluant directement la qualité de la décision d'inspection.
\end{frame}

\begin{frame}{Comment interpréter les métriques du détecteur ?}
\begin{itemize}
  \item \textbf{Précision} : parmi les alertes, quelle part correspond à un vrai défaut ?
  \item \textbf{Rappel} : parmi les vrais défauts, quelle part est détectée ?
  \item \textbf{F1-score} : compromis entre précision et rappel.
  \item \textbf{ROC AUC / PR AUC} : qualité globale de séparation sur différents seuils.
\end{itemize}

\vspace{0.2cm}
Cette lecture est pertinente car une bonne politique conditionnelle dépend fortement de la qualité du détecteur.
\end{frame}

\section{Validation et implémentation}

\begin{frame}{Validation du simulateur}
\begin{itemize}
  \item Des tests automatiques ont été écrits pour vérifier le bon fonctionnement du simulateur.
  \item Les points vérifiés incluent :
  \begin{itemize}
    \item couverture complète du journal quotidien,
    \item bon fonctionnement des lois d'usure,
    \item déclenchement des maintenances,
    \item génération correcte des sorties Monte Carlo.
  \end{itemize}
  \item Les exports CSV facilitent aussi la traçabilité du comportement simulé.
\end{itemize}
\end{frame}

\begin{frame}{Bonus : inversion du problème}
\begin{itemize}
  \item Le problème direct consiste à fixer les paramètres puis simuler les pannes.
  \item Le problème inverse consiste à partir de données observées pour retrouver les paramètres du modèle.
  \item Nous avons abordé ce bonus via une calibration simple par simulation :
  \begin{itemize}
    \item estimation directe de la loi d'usage à partir des usages observés ;
    \item recherche de paramètres de dégradation et de panne capables de reproduire les statistiques observées.
  \end{itemize}
  \item Cette étape est exploratoire, mais elle montre que le simulateur peut servir de base à une future calibration sur données réelles.
\end{itemize}
\end{frame}

\begin{frame}{Bonus : qualité de la calibration inverse}
\small
\centering
\begin{tabular}{lrrr}
\toprule
\textbf{Métrique} & \textbf{Observé} & \textbf{Simulé} & \textbf{Erreur rel.} \\
\midrule
Nb. pannes & 58 & 57 & 1.7\% \\
Nb. inspections & 4742 & 4732 & 0.2\% \\
Coût total & 1671285 & 1880534 & 12.5\% \\
Disponibilité & 0.9776 & 0.9722 & 0.5\% \\
Score insp. moyen & 0.8118 & 0.8275 & 1.9\% \\
\bottomrule
\end{tabular}

\vspace{0.3cm}
\begin{itemize}
  \item Les données observées utilisées ici proviennent de la politique \texttt{conditionnelle\_30j}.
  \item La calibration reproduit correctement plusieurs grandeurs globales, mais certaines statistiques structurelles restent plus difficiles à ajuster.
\end{itemize}
\end{frame}

\begin{frame}{Organisation des sorties}
\begin{itemize}
  \item \texttt{comparaison\_politiques.csv} : synthèse globale
  \item un dossier par politique avec :
  \begin{itemize}
    \item \texttt{usage\_log.csv}
    \item \texttt{inspections.csv}
    \item \texttt{maintenance.csv}
    \item \texttt{pannes.csv}
    \item \texttt{etat\_journalier.csv}
  \end{itemize}
  \item \texttt{monte\_carlo/} : robustesse statistique
  \item \texttt{ml\_metrics/} : scores de classification et courbes ROC / PR
\end{itemize}
\end{frame}

\section{Conclusion}

\begin{frame}{Conclusion, limites et perspectives}
\begin{itemize}
  \item Le simulateur permet de comparer plusieurs politiques de maintenance dans un cadre probabiliste cohérent.
  \item La maintenance corrective pure est clairement défavorable dans notre paramétrage.
  \item Les politiques conditionnelle et systématique constituent les meilleurs compromis observés.
\end{itemize}

\vspace{0.3cm}
\textbf{Limites :}
\begin{itemize}
  \item paramètres non calibrés sur données réelles,
  \item modélisation mono-composant par système,
  \item pas de contrainte de stock ni de ressource équipe.
\end{itemize}

\textbf{Perspectives :}
\begin{itemize}
  \item calibrer les paramètres sur des données observées réelles,
  \item étendre le modèle à plusieurs composants avec interactions,
  \item intégrer les contraintes opérationnelles : stock, équipes, délais et optimisation de politique.
\end{itemize}
\end{frame}

\begin{frame}{Message final}
\begin{block}{Conclusion du projet}
Nous avons construit un simulateur stochastique simple, lisible et exploitable pour comparer des politiques de maintenance sur une flotte de systèmes.
\end{block}

\vspace{0.25cm}
\begin{itemize}
  \item modèle probabiliste cohérent ;
  \item comparaison économique et opérationnelle ;
  \item quantification de l'incertitude ;
  \item lecture complémentaire via des métriques de classification.
\end{itemize}

\vspace{0.25cm}
\centering
\Large Merci pour votre attention.
\end{frame}

\end{document}
